\documentclass[11pt]{article}
\usepackage[margin=1in]{geometry}
\usepackage[T1]{fontenc}
\usepackage[utf8]{inputenc}
\usepackage{amsmath,amssymb,mathtools}
\usepackage{booktabs}
\usepackage{hyperref}
\usepackage{xcolor}
\usepackage{listings}

\hypersetup{
  colorlinks=true,
  linkcolor=blue!60!black,
  urlcolor=blue!60!black
}

\definecolor{codebg}{RGB}{248,248,248}
\definecolor{codeborder}{RGB}{210,210,210}
\definecolor{codecomment}{RGB}{80,120,80}
\definecolor{codekeyword}{RGB}{30,60,150}

\lstdefinestyle{reportcode}{
  backgroundcolor=\color{codebg},
  basicstyle=\ttfamily\small,
  breaklines=true,
  columns=fullflexible,
  commentstyle=\color{codecomment},
  frame=single,
  framerule=0.5pt,
  keywordstyle=\color{codekeyword},
  rulecolor=\color{codeborder},
  keepspaces=true,
  showstringspaces=false,
  tabsize=2
}

\lstset{style=reportcode, language=Python}

\setlength{\parindent}{0pt}
\setlength{\parskip}{0.5em}

\newcommand{\R}{\mathbb{R}}
\newcommand{\ysp}{y_{\mathrm{sp}}}
\newcommand{\uvec}{u}
\newcommand{\xvec}{x}
\newcommand{\Aaug}{A_{\mathrm{aug}}}
\newcommand{\Baug}{B_{\mathrm{aug}}}
\newcommand{\Caug}{C_{\mathrm{aug}}}
\newcommand{\Qout}{Q_{\mathrm{out}}}
\newcommand{\Rin}{R_{\mathrm{in}}}
\newcommand{\Hp}{H_p}
\newcommand{\Hc}{H_c}
\newcommand{\argmin}{\mathop{\mathrm{arg\,min}}}


\title{Weight Adaptation In RL-Assisted MPC}
\date{2026-04-01}

\begin{document}
\maketitle

\section*{Method Goal}

This method uses a TD3 agent to adapt the four scalar penalties that define the
MPC stage cost: two output-tracking weights and two input-move weights. The
continuous policy does not replace the MPC optimizer. It changes the optimizer's
relative priorities online by multiplying the baseline $\Qout$ and $\Rin$
entries.

\section*{Notebook Sources Covered}

\begin{itemize}
\item \texttt{RL\_assisted\_MPC\_weights.ipynb}
\item \texttt{RL\_assisted\_MPC\_weights\_dist.ipynb}
\item Shared helpers from \texttt{TD3Agent/agent.py} and
      \texttt{Simulation/mpc.py}
\end{itemize}

\section*{Shared MPC Context}

The baseline MPC stays fixed in structure:

\begin{itemize}
\item the plant still receives the first move of an optimization-based MPC,
\item the observer state still drives the internal prediction model,
\item the RL policy only changes how expensive different errors and moves are.
\end{itemize}

If the baseline penalties are
\[
\Qout^{(0)} = [Q_1, Q_2],
\qquad
\Rin^{(0)} = [R_1, R_2],
\]
then the agent chooses a four-vector of multipliers
\[
m_t = [m_{Q_1}, m_{Q_2}, m_{R_1}, m_{R_2}]
\]
and the active MPC penalties become
\[
\Qout^{(t)} = [Q_1 m_{Q_1}, Q_2 m_{Q_2}],
\qquad
\Rin^{(t)} = [R_1 m_{R_1}, R_2 m_{R_2}].
\]

\section*{State, Action, Reward, And Update Logic}

The state is the same scaled observer-state bundle used by most TD3 notebook
families:
\[
s_t =
\begin{bmatrix}
\tilde{x}_{t} \\
\tilde{y}_{\mathrm{sp},t} \\
\tilde{u}_{t}^{\mathrm{dev}}
\end{bmatrix}.
\]

The action is continuous with dimension four. A raw TD3 action
$a_t \in [-1, 1]^4$ is remapped into user-defined bounds:
\[
m_t = \ell + \frac{a_t + 1}{2} \odot (h - \ell),
\]
where $\ell$ and $h$ are the lower and upper coefficient bounds.

The reward still comes from closed-loop tracking and move quality after solving
the weighted MPC problem. This notebook family also multiplies the reward by
$0.01$ before storing it.

\section*{Core Algorithm Flow}

\begin{enumerate}
\item Build the scaled RL state from the observer state and operating point.
\item Query the TD3 agent for a four-dimensional action.
\item Map that action to four penalty multipliers.
\item Update $\Qout$ and $\Rin$ inside the current MPC object.
\item Solve MPC, apply the first move to the plant, and roll the observer.
\item Store the transition in the TD3 replay buffer and call
      \texttt{train\_step()} after warm start.
\end{enumerate}

\section*{Key Code Paths}

\subsection*{Notebook-local penalty remapping}

\begin{lstlisting}
Q_base = np.array([Q1_penalty, Q2_penalty], dtype=float)
R_base = np.array([R1_penalty, R2_penalty], dtype=float)

def map_to_bounds(a, low, high):
    return low + ((a + 1.0) / 2.0) * (high - low)

def set_penalties(MPC_obj, Q_base, R_base, m):
    Q1_new = float(Q_base[0] * m[0])
    Q2_new = float(Q_base[1] * m[1])
    R1_new = float(R_base[0] * m[2])
    R2_new = float(R_base[1] * m[3])
    MPC_obj.Q_out = np.array([Q1_new, Q2_new])
    MPC_obj.R_in = np.array([R1_new, R2_new])
\end{lstlisting}

\subsection*{Notebook-local TD3 control loop}

\begin{lstlisting}
if i > WARM_START:
    if not test:
        action = agent.take_action(current_rl_state, explore=(not test))
    else:
        action = agent.act_eval(current_rl_state)
else:
    action_ = np.array([1., 1., 1., 1.])
    action = map_from_bounds(action_, low_coef, high_coef)

m = map_to_bounds(action, low_coef, high_coef)
set_penalties(MPC_obj, Q_base, R_base, m)

sol = spo.minimize(
    lambda x: MPC_obj.mpc_opt_fun(x, y_sp[i, :],
                                  scaled_current_input_dev,
                                  xhatdhat[:, i]),
    IC_opt, bounds=bnds, constraints=cons
)
\end{lstlisting}

\subsection*{Shared TD3 interaction and update logic}

\begin{lstlisting}
@torch.no_grad()
def take_action(self, state, explore=False):
    self.steps += 1
    s = torch.as_tensor(state, dtype=torch.float32, device=self.device)
    a = self.actor(s).detach().cpu().numpy()
    if explore:
        self._expl_sigma = self.expl_sched.value(self.steps)
        a = a + np.random.randn(*a.shape) * self._expl_sigma
    return np.clip(a, -self.max_action, self.max_action)

def train_step(self):
    ...
    next_action = (base_next + noise).clip(-self.max_action, self.max_action)
    q_next = self.critic_target.combined_forward(ns, next_action,
                                                 mode=self.target_combine)
    y = r + self.gamma * (1.0 - done) * q_next
\end{lstlisting}

\section*{Variants Covered}

\begin{itemize}
\item The nominal notebook adapts weights without injected process drift.
\item The disturbance notebook applies gradual changes to process quantities
      while reusing the same TD3-based weight tuning loop.
\item Both notebooks log the active weight multipliers over time through
      \texttt{mult\_log}.
\end{itemize}

\section*{Limitations, Caveats, And Open Questions}

\begin{itemize}
\item This family assumes the MPC structure is correct and only tunes the cost.
\item The weight bounds are notebook-level configuration, not a shared schema.
\item The control loop is still notebook-local, so refactors must preserve the
      notebook supervisor behavior instead of editing only module helpers.
\item Historical comparisons depend on preserving the same baseline
      $Q$ and $R$ definitions.
\end{itemize}

\end{document}
